\section*{The Propers: Themes and Movement}

\begin{massbox}{The newly alive remain a dependent people}
Romans announces a real baptismal passage into Christ's death and life; the surrounding prayers refuse to treat that gift as self-sufficiency. God strengthens, implants, nourishes, and guards; Christ feeds bodies that would fail on the road; the offerer asks for steadied steps and faithful prayer; Communion turns received mercy into jubilation while the Postcommunion still asks for cleansing and aid.
\end{massbox}

\begin{center}
\begin{tikzpicture}[x=3.55cm]
\node[guide node,text width=2.65cm] (font) at (0,0) {\textbf{Font}\\buried, planted, raised};
\node[guide node,text width=2.65cm] (desert) at (1,0) {\textbf{Desert}\\refuge and broken bread};
\node[guide emphasis,text width=2.65cm] (road) at (2,0) {\textbf{Road}\\inclined ear and steady steps};
\node[guide node,text width=2.65cm] (altar) at (3,0) {\textbf{Altar}\\effective prayer and fortifying gift};
\foreach \x/\y in {font/desert,desert/road,road/altar}
  \draw[guide arrow] (\x) -- (\y);
\end{tikzpicture}
\end{center}

\noindent The itinerary runs from font to altar: receive life at the font; seek refuge and receive food in the desert; ask for hearing and steady steps along the road; offer and receive at the altar. The Fathers cited below expound the separate texts; the itinerary is a reading of their sequence.

\sectionguard
\subsection*{Strength Gives a New Walk\enspace\properrefs{Int., Coll., Ep.}}

The Introit joins a cry against divine silence to confession of the Lord as his people's strength, his anointed's saving protector, and the one who blesses and shepherds the inheritance. The Collect keeps that agency before every later command: God implants love of his name, increases religion, nourishes the good, and guards what he has nourished. Paul then states the baptismal indicative before the moral imperative. Those buried with Christ are commanded to walk and reckon according to a life they received rather than produced.

\begin{studybox}{Chrysostom on Romans 6}
Chrysostom distinguishes Christ's bodily death, the baptized person's real death to sin, continuing mortification with divine help, present fruit, and future bodily resurrection. The body is not the ``old man'' and is not evil. Grace establishes a new reality that summons cooperation without becoming a merely honorary status.
\end{studybox}

\sectionguard
\subsection*{Refuge and Compassion Meet Mortal Need\enspace\properrefs{Grad., Gosp.}}

The Moses-titled Gradual places mortal generations beneath the God who has always been their refuge, then asks him to return and show favor to his servants. Mark shows that dependence bodily: a concrete multitude has stayed with Jesus for three days and would collapse if dismissed fasting. Christ's human compassion and divine power preserve bodily need; Bede and Augustine further read the broken bread as opened Scripture and ecclesial nourishment without erasing the literal meal.

\begin{threestable}{Witness}{Evidence}{Distinction}
Psalm 89 & God's eternity and generational refuge frame mortality, wrath, and the plea of servants for his return. & The communal lament has its own scriptural horizon, prior to later Christian reception.\\
Mark 8 & Christ pities and feeds four thousand people who cannot sustain themselves on the journey. & Bodies, hunger, loaves, fish, and material mercy remain literal.\\
Bede and Augustine & Broken bread opens scriptural mysteries; apostolic ministry distributes the food of life, and seven can signify the Spirit's complete work. & Their direct exegesis is chiefly scriptural and ecclesial, not a Eucharistic code for every detail.\\
\end{threestable}

\sectionguard
\subsection*{One Inclined Ear Holds Hope and Steps\enspace\properrefs{All., Off.}}

The Alleluia and Offertory repeat the same core appeal, \latin{inclina aurem tuam}. Confidence therefore sounds as petition rather than self-possession, and steadfastness remains something God is asked to secure.

\begin{threestable}{Proper}{Textual action}{What the echo establishes}
Alleluia & \latin{In te ... speravi}; deliver and rescue; \latin{inclina ad me aurem tuam}; hasten to rescue. & Before the Gospel, threatened trust places danger and possible shame beneath God's righteousness and hearing.\\
Offertory & Perfect steps so footprints do not move; \latin{inclina aurem tuam}; hear; make mercies wonderful for those who hope. & At the offering, the same bent ear is sought for steadfast conduct and spoken petition.\\
\end{threestable}

\sectionguard
\subsection*{Faithful Asking Becomes Jubilation and Aid\enspace\properrefs{Sec., Comm., Postcomm.}}

The three altar texts refuse to turn offering, praise, or reception into self-sufficiency. Their sequence moves from gifts presented, through jubilation at Communion, to effects still requested after fullness.

\begin{threestable}{Proper}{Evidence}{Distinction}
Secret & God is asked to receive the people's oblations and grant that faithful petition attain \latin{efficaciter}. & Augustine's \emph{Letter 130} prevents efficacy from becoming a technique for compelling a preferred temporal outcome.\\
Communion & Psalm 26 offers a sacrifice of jubilation; Augustine reads it as Christ-centered confession sounding in heart and deed. & Augustine directly expounds the psalm, not its later Communion placement or an explicit Eucharistic meaning.\\
Postcommunion & Those filled with God's gifts ask through paired passives to be cleansed by their effect and fortified by their aid. & Fullness is real, but reception does not end the gifts' operation or the recipients' dependence.\\
\end{threestable}

\begin{studybox}{Three orders held together}
\textbf{Indicative:} God has joined the baptized to Christ. \textbf{Imperative:} therefore walk and reckon according to the gift. \textbf{Petition:} ask God to implant, nourish, hear, fulfill, cleanse, and fortify. Removing the indicative produces moralism; removing the imperative produces inert sacramentalism; removing petition imagines that grace is needed only at the beginning.
\end{studybox}
